\documentclass{article}
\usepackage{nips14submit_e,times}
%%%%\usepackage{hyperref}
\usepackage{url}

\usepackage{graphicx}
\usepackage{subfigure}
\usepackage{framed}
\usepackage{wrapfig}
\usepackage{tabularx}

% For citations
\usepackage[numbers]{natbib}

% For algorithms
\usepackage{algorithm}
\usepackage{algorithmic}

\usepackage{amsmath}
\usepackage{amsthm}
\usepackage{amssymb}
\usepackage{pifont}
\usepackage{rotating}

\newtheorem{thm}{Theorem}
\newtheorem{lem}{Lemma}
\newtheorem{cor}{Corollary}

\setlength{\parskip}{5pt}
\setlength{\abovedisplayskip}{4pt plus2pt minus3pt}
\setlength{\belowdisplayskip}{\abovedisplayskip}

\setlength{\abovedisplayshortskip}{ 2pt plus3pt}
\setlength{\belowdisplayshortskip}{4pt plus3pt minus3pt}

\allowdisplaybreaks[4]
\newcommand{\cmark}{\ding{51}}
\newcommand{\xmark}{\ding{55}}

\newcommand{\argmin}{\operatornamewithlimits{argmin}}
\newcommand{\argmax}{\operatornamewithlimits{argmax}}

\newcommand{\stepsize}{\gamma} % more distinguishable from n than eta!



 



%%%%%%%%%%%%%%%
% Default fixed font does not support bold face
\DeclareFixedFont{\ttb}{T1}{txtt}{bx}{n}{8} % for bold
\DeclareFixedFont{\ttm}{T1}{txtt}{m}{n}{8}  % for normal

% Custom colors
\usepackage{color}
\definecolor{deepblue}{rgb}{0,0,0.5}
\definecolor{deepred}{rgb}{0.6,0,0}
\definecolor{deepgreen}{rgb}{0,0.5,0}

\usepackage{listings}

% Python style for highlighting
\newcommand\pythonstyle{\lstset{
language=Python,
basicstyle=\ttm,
otherkeywords={Self},             % Add keywords here
keywordstyle=\ttb\color{deepblue},
emph={cdef,double,long,int,unsigned,inline},          % Custom highlighting
emphstyle=\ttb\color{deepred},    % Custom highlighting style
stringstyle=\color{deepgreen},
commentstyle=\color{deepgreen},
frame=tb,                         % Any extra options here
showstringspaces=false            % 
}}


% Python environment
\lstnewenvironment{python}[1][]
{
\pythonstyle
%\lstset{#1}
}
{}



\title{SAGA: A Fast Incremental Gradient Method With Support for Non-Strongly Convex Composite Objectives}

\author{
Aaron Defazio \\
Ambiata \thanks{The first author completed this work while under funding from NICTA. This work was partially supported by the MSR-Inria Joint Centre and a grant by the European Research Council (SIERRA project 239993).} \\
Australian National University, Canberra\\
\And
Francis Bach \\
INRIA - Sierra Project-Team \\ 
\'Ecole Normale Sup\'erieure, Paris, France \\
\AND
Simon Lacoste-Julien \\
INRIA - Sierra Project-Team \\
\'Ecole Normale Sup\'erieure, Paris, France \\
}

\newcommand{\fix}{\marginpar{FIX}}
\newcommand{\new}{\marginpar{NEW}}


\def \E{{\mathbb E}} % expectation operator
\def \V { \mathbb{V} } % variance operator
\DeclareMathOperator{\var}{Var} % variance
\DeclareMathOperator{\cov}{Cov} % variance



\nipsfinalcopy

\begin{document}

\maketitle


\begin{abstract}
\vspace{-2mm}
In this work we introduce a new optimisation method called SAGA in the spirit of SAG, SDCA, MISO and SVRG, a set of recently proposed incremental gradient algorithms with fast linear convergence rates. SAGA improves on the theory behind SAG and SVRG, with better theoretical convergence rates, and has support for composite objectives where a proximal operator is used on the regulariser. Unlike SDCA, SAGA supports non-strongly convex problems directly, and is adaptive to any inherent strong convexity of the problem. We give experimental results showing the effectiveness of our method.
\end{abstract}

\section{Introduction}
\label{sec:intro}
\vspace{-2mm}
Remarkably, recent advances \citep{SAG, SDCA} have shown that it is possible to minimise strongly convex finite sums provably faster in expectation than is possible without the finite sum structure. This is significant for machine learning problems as a finite sum structure is common in the empirical risk minimisation setting. The requirement of strong convexity is  likewise satisfied in machine learning problems in the typical case where a quadratic regulariser is used.  

In particular, we are interested in minimising functions of the form
\[
f(x)=\frac{1}{n}\sum_{i=1}^{n}f_{i}(x), 
\]
where $x\in\mathbb{R}^d$, each $f_{i}$ is convex and has Lipschitz continuous derivatives with constant $L$.
We will also consider the case where each $f_{i}$ is strongly convex with constant $\mu$, and the ``composite" (or proximal) case where an additional regularisation function is added:
\[
F(x) = f(x) + h(x),
\]
where $h\colon \mathbb{R}^d \rightarrow \mathbb{R}^d$ is convex but potentially non-differentiable, and where the proximal operation of $h$ is easy to compute --- few incremental gradient methods are applicable in this setting~\citep{jota}\citep{prox-svrg}. %SLJ: rephrase to save one line...

Our contributions are as follows. In Section \ref{sec:algorithm} we describe the SAGA algorithm, a novel incremental gradient method. In Section \ref{sec:theory} we prove theoretical convergence rates for SAGA in the strongly convex case better than those for SAG \citep{SAG} and SVRG \citep{svrg}, and a factor of 2 from the SDCA \citep{SDCA} convergence rates. These rates also hold in the composite setting. Additionally, we show that like SAG but unlike SDCA, our method is applicable to non-strongly convex problems without modification. We establish theoretical convergence rates for this case also.
In Section \ref{sec:related-work} we discuss the relation between each of the fast incremental gradient methods, showing that each stems from a very small modification of another.
\vspace{-2mm}
\section{SAGA Algorithm}
\label{sec:algorithm}
\vspace{-2mm}
We start with some known initial vector $x^{0}\in\mathbb{R}^d$ and known derivatives $f_{i}^{\prime}(\phi^0_i)\in\mathbb{R}^d$ with $\phi^0_i = x^{0}$ for
each $i$. These derivatives are stored in a table data-structure of length $n$, or alternatively a $n\times d$ matrix.
For many problems of interest, such as binary classification and least-squares, only a single floating point value instead of a full gradient vector needs to be stored (see Section~\ref{sec:impl}).
SAGA is inspired both from SAG~\citep{SAG} and SVRG~\citep{svrg} (as we will discuss in Section~\ref{sec:related-work}). SAGA uses a step size of $\stepsize$ and makes the following updates, starting with $k=0$:

% adefazio: This was too wide before, so it didn't pass NIPS validation.
\noindent\framebox{\begin{minipage}[t]{0.98\textwidth}
\textbf{SAGA Algorithm:} Given the value of $x^{k}$ and of each $f_{i}^{\prime}(\phi_{i}^{k})$
at the end of iteration $k$, the updates for iteration $k+1$ is
as follows:
\begin{enumerate}
\item Pick a $j$ uniformly at random.
\item Take $\phi_{j}^{k+1}=x^{k}$, and store $f_{j}^{\prime}(\phi_{j}^{k+1})$ in the table. All other entries in the table remain unchanged. The quantity $\phi_{j}^{k+1}$ is not explicitly stored.
\item Update $x$ using $f_{j}^{\prime}(\phi_{j}^{k+1})$, $f_{j}^{\prime}(\phi_{j}^{k})$ and the table average:
\begin{equation}
w^{k+1}=x^{k}-\stepsize  \left[ f_{j}^{\prime}(\phi_{j}^{k+1})-f_{j}^{\prime}(\phi_{j}^{k})+\frac{1}{n}\sum^{n}_{i=1}f_{i}^{\prime}(\phi_{i}^{k})\right], \label{eq:main-update}
\end{equation}
\begin{equation}
x^{k+1}=\text{prox}^{h}_{\stepsize}\left(w^{k+1}\right).
\end{equation}
\end{enumerate}
\end{minipage}}

The proximal operator we use above is defined as
\begin{equation}
\text{prox}^{h}_{\stepsize}\left(y\right) := \argmin_{x \in \mathbb{R}^d} \,\left\{ h(x) + \frac{1}{2\stepsize} \Vert x - y \Vert^2 \right\}.
\end{equation}
In the strongly convex case, when a step size of $\stepsize=1/(2(\mu n+L))$ is chosen, we have the following convergence rate in the composite and hence also the non-composite case:
\begingroup\makeatletter\def\f@size{9}\check@mathfonts
\[
\mathbb{E}\left\Vert x^{k}-x^{*}\right\Vert ^{2}\leq\left(1-\frac{\mu}{2(\mu n+L)}\right)^{k}\left[\left\Vert x^{0}-x^{*}\right\Vert ^{2}
+\frac{n}{\mu n + L} \left[f(x^{0}) -\left\langle f^{\prime}(x^{*}),x^{0}-x^{*}\right\rangle  - f(x^{*})\right]\right].
 \]
 \endgroup
We prove this result in Section \ref{sec:theory}. The requirement of strong convexity can be relaxed from needing to hold for each $f_i$ to just holding on average, but at the expense of a worse geometric rate ($1-\frac{\mu}{6(\mu n+L)})$, requiring a step size of $\stepsize=1/(3(\mu n+L))$.

In the non-strongly convex case, we have established the convergence rate in terms of the average iterate, excluding step 0:  $\bar{x}^{k}=\frac{1}{k}\sum_{t=1}^{k}x^{t}$. Using a step size of $\stepsize=1/(3L)$ we have
\[
\mathbb{E}\left[F(\bar{x}^{k})\right]-F(x^{*})\leq\frac{4n}{k}\left[\frac{2L}{n}\left\Vert x^{0}-x^{*}\right\Vert ^{2}
+f(x^{0}) -\left\langle f^{\prime}(x^{*}),x^{0}-x^{*}\right\rangle  -f(x^{*})\right].
\]
This result is proved in the supplementary material. Importantly, when this   step size $\stepsize=1/(3L)$ is used, our algorithm \emph{automatically adapts} to the level of strong convexity $\mu > 0$ naturally present, giving a convergence rate of (see the comment at the end of the proof of Theorem~\ref{thm:strong-convex-thm}):
\begingroup\makeatletter\def\f@size{9}\check@mathfonts
\[
\mathbb{E}\left\Vert x^{k}-x^{*}\right\Vert ^{2}\leq
\left( 1-\min\left\{ \frac{1}{4n}\,,\,\frac{\mu}{3L}\right\} \right)^{k}
\left[\left\Vert x^{0}-x^{*}\right\Vert ^{2}
+\frac{2n}{3L} \left[f(x^{0}) -\left\langle f^{\prime}(x^{*}),x^{0}-x^{*}\right\rangle  -f(x^{*})\right]\right].
 \]
 \endgroup
Although any incremental gradient method can be applied to non-strongly convex problems via the addition of a small quadratic regularisation, the amount of regularisation is an additional tunable parameter which our method avoids.
 
\vspace{-2mm}
\section{Related Work}
\label{sec:related-work}
\vspace{-2mm}
We explore the relationship between SAGA and the other fast incremental gradient methods in this section. By using SAGA as a midpoint, we are able to provide a more unified view than is available in the existing literature. A brief summary of the properties of each method considered in this section is given in Figure~\ref{fig:comparison-table}. The method from~\citep{jota}, which handles the non-composite setting, is not listed as its rate is of the slow type and can be up to $n$ times smaller than the one for SAGA or SVRG~\citep{svrg}.

\begin{figure}[t]
\center
\small
\begin{tabular*}{\textwidth}{@{\extracolsep{\fill}} |cccccc|}
\hline 
 & SAGA & SAG & SDCA & SVRG & FINITO \tabularnewline
Strongly Convex (SC) & \cmark & \cmark & \cmark & \cmark & \cmark \tabularnewline
Convex, Non-SC* & \cmark & \cmark & \xmark & ? & ? \tabularnewline
Prox Reg. & \cmark & ? & \cmark \cite{sdca-admm} & \cmark & \xmark \tabularnewline 
Non-smooth & \xmark & \xmark & \cmark & \xmark & \xmark \tabularnewline
Low Storage Cost & \xmark & \xmark & \xmark & \cmark & \xmark \tabularnewline
Simple(-ish) Proof & \cmark & \xmark & \cmark & \cmark & \cmark \tabularnewline 
Adaptive to SC & \cmark & \cmark & \xmark & ? & ? \tabularnewline
\hline
\end{tabular*}
\vskip -0.1in
\caption{\small \label{fig:comparison-table} Basic summary of method properties. Question marks denote unproven, but not experimentally ruled out cases. (*) Note that any method can be applied to non-strongly convex problems by adding a small amount of L2 regularisation, this row describes methods that do not require this trick. \vspace{-3mm}}
\end{figure}


% --- old subsetion SAG with u formulation --- 
%\subsection*{SAG}
%\label{subsec:sag}
%If we eliminate $x^{k}$, we get an update for $u$ in SAGA of:
%\begin{equation}
%u^{k+1}=u^{k}-\frac{\stepsize}{n}\sum^{n}_{i=1}f_{i}^{\prime}(\phi_{i}^{k}). \label{eq:u-update}
%\end{equation}
%After translating notation, this is identical to the SAG (Stochastic Average Gradient) \citep{SAG} update, except instead of setting $\phi_{j}^{k+1}=u^{k}$, we are using an update of $\phi_{j}^{k+1}=x^{k}=u^{k}-\stepsize \sum_{i}f_{i}^{\prime}(\phi_{i}^{k})$.
%The order of the steps is also changed, as in SAG the $j$th gradient
%is updated before the $x$ step is taken, where as above it is updated
%after. However the order of the steps for SAG doesn't affect the algorithm so
%the ordering change is not significant.
%
%The significance of this change is the effect it has on our current estimate of the gradient. In SAG, the gradient approximation $\frac{1}{n}\sum_{i}f_{i}^{\prime}(\phi_{i}^{k})$ is biased away from the true gradient, whereas for all the other methods considered here, including SAGA, the gradient approximation is unbiased. The trade-off for the increased bias is a decreased variance in the $f_{i}^{\prime}(\phi_{i}^{k})$ gradients, due to the different $\phi_{i}^{k}$ updates used.
%
%The per update cost of SAGA and SAG is essentially the same. The advantage over SAG is that SAGA has a much more complete, simple and tight theory. The theoretical convergence rate is better for SAGA, and no theory exists for the use of proximal operators in SAG.

\vspace{-2mm}
\subsection*{SAGA: midpoint between SAG and SVRG/S2GD}
\label{subsec:sag}

In~\citep{svrg}, the authors make the observation that the variance of the standard stochastic gradient (SGD) update direction can only go to zero if decreasing step sizes are used, thus preventing a linear convergence rate unlike for batch gradient descent. They thus propose to use a variance reduction approach (see~\citep{variance-reduction} and references therein for example) on the SGD update in order to be able to use constant step sizes and get a linear convergence rate. We present the updates of their method called SVRG (Stochastic Variance Reduced Gradient) in~\eqref{eq-SVRG} below, comparing it with the non-composite form of SAGA rewritten in~\eqref{eq:w-update}. They also mention that SAG (Stochastic Average Gradient)~\citep{SAG} can be interpreted as reducing the variance, though they do not provide the specifics. Here, we make this connection clearer and relate it to SAGA. 

We first review a slightly more generalized version of the variance reduction approach (we allow the updates to be biased). Suppose that we want to use Monte Carlo samples to estimate $\E X$ and that we can compute efficiently $\E Y$ for another random variable $Y$ that is highly correlated with $X$. One variance reduction approach is to use the following estimator $\theta_\alpha$ as an approximation to $\E X$: $\theta_\alpha := \alpha(X-Y) + \E Y$, for a step size $\alpha \in [0,1]$. We have that $\E \theta_\alpha$ is a convex combination of $\E X$ and $\E Y$: $\E \theta_\alpha = \alpha \E X + (1-\alpha) \E Y$. The standard variance reduction approach uses $\alpha=1$ and the estimate is unbiased $\E \theta_1 = \E X$. The variance of $\theta_\alpha$ is: $\var(\theta_\alpha) = \alpha^2 [\var(X) + \var(Y) - 2 \cov(X,Y)]$, and so if $\cov(X,Y)$ is big enough, the variance of $\theta_\alpha$ is reduced compared to $X$, giving the method its name. By varying $\alpha$ from 0 to 1, we increase the variance of $\theta_\alpha$ towards its maximum value (which usually is still smaller than the one for $X$) while decreasing its bias towards zero. 

Both SAGA and SAG can be derived from such a variance reduction viewpoint: here $X$ is the SGD direction sample $f_{j}^{\prime}(x^k)$, whereas $Y$ is a past stored gradient $f_{j}^{\prime}(\phi_j^k)$. SAG is obtained by using $\alpha=1/n$ (update rewritten in our notation in~\eqref{SAG}), whereas SAGA is the unbiased version with $\alpha=1$ (see~\eqref{eq:w-update} below). For the same $\phi$'s, the variance of the SAG update is $1/n^2$ times the one of SAGA, but at the expense of having a non-zero bias. This non-zero bias might explain the complexity of the convergence proof of SAG and why the theory has not yet been extended to proximal operators. By using an unbiased update in SAGA, we are able to obtain a simple and tight theory, with better constants than SAG, as well as theoretical rates for the use of proximal operators.

\vspace{-5mm}
\begin{align}
\textrm{(SAG)} \qquad x^{k+1} &=x^{k}-\stepsize  \left[ \frac{f_{j}^{\prime}(x^k)-f_{j}^{\prime}(\phi_{j}^{k})}{n}+\frac{1}{n}\sum^{n}_{i=1}f_{i}^{\prime}(\phi_{i}^{k})\right], \label{SAG} \\
\textrm{(SAGA)} \qquad x^{k+1} &=x^{k}-\stepsize  \left[ f_{j}^{\prime}(x^k)-f_{j}^{\prime}(\phi_{j}^{k})+\frac{1}{n}\sum^{n}_{i=1}f_{i}^{\prime}(\phi_{i}^{k})\right], \label{eq:w-update} \\
\textrm{(SVRG)} \qquad x^{k+1} &=x^{k}-\stepsize  \left[ f_{j}^{\prime}(x^{k})-f_{j}^{\prime}(\tilde{x})+\frac{1}{n}\sum^{n}_{i=1}f_{i}^{\prime}(\tilde{x})\right]. \label{eq-SVRG}
\end{align}

The SVRG update~\eqref{eq-SVRG} is obtained by using $Y=f_{j}^{\prime}(\tilde{x})$ with $\alpha=1$ (and is thus unbiased -- we note that SAG is the only method that we present in the related work that has a biased update direction).
The vector $\tilde{x}$ is not updated every step, but rather the loop over $k$ appears inside an outer loop, where $\tilde{x}$ is updated at the start of each outer iteration. Essentially SAGA is at the midpoint between SVRG and SAG; it updates the $\phi_{j}$ value each time index $j$ is picked, whereas SVRG updates all of $\phi$'s as a batch. 
The S2GD method \citep{semi} has the same update as SVRG, just differing in how the number of inner loop iterations is chosen. We use SVRG henceforth to refer to both methods.

SVRG makes a trade-off between time and space. For the equivalent practical convergence rate it makes 2x-3x more gradient evaluations, but in doing so it does not need to store a table of gradients, but a single average gradient. The usage of SAG vs. SVRG is problem dependent. For example for linear predictors where gradients can be stored as a reduced vector of dimension $p-1$ for $p$ classes, SAGA  is preferred over SVRG both theoretically and in practice. For neural networks, where no theory is available for either method, the storage of gradients is generally more expensive than the additional backpropagations, but this is computer architecture dependent. 

SVRG also has an additional parameter besides step size that needs to be set, namely the number of iterations per inner loop ($m$). This parameter can be set via the theory, or conservatively as $m=n$, however doing so does not give anywhere near the best practical performance. Having to tune one parameter instead of two is a practical advantage for SAGA.

\subsection*{Finito/MISO$\mu$ }
\label{subsec:finito}
To make the relationship with other prior methods more apparent, we can rewrite the SAGA algorithm (in the non-composite case) in term of an additional intermediate quantity $u^{k}$, with $u^{0} := x^{0}+ \stepsize \sum_{i=1}^{n}f_{i}^{\prime}(x^{0})$, in addition to the usual $x^{k}$ iterate as described previously: 

\noindent\framebox{\begin{minipage}[t]{0.98\textwidth}
\textbf{SAGA: Equivalent reformulation for non-composite case:}
Given the value of $u^{k}$ and of each $f_{i}^{\prime}(\phi_{i}^{k})$
at the end of iteration $k$, the updates for iteration $k+1$, is
as follows:
\begin{enumerate}
\item Calculate $x^k$:
\vspace{-4mm}
\begin{equation}
x^{k}=u^{k}-\stepsize \sum_{i=1}^{n}f_{i}^{\prime}(\phi_{i}^{k}). \label{eq:alt-form-update}
\end{equation}

\item Update $u$ with $u^{k+1}=u^{k}+\frac{1}{n}(x^{k}-u^{k})$.
\item Pick a $j$ uniformly at random.
\item Take $\phi_{j}^{k+1}=x^{k}$, and store $f_{j}^{\prime}(\phi_{j}^{k+1})$ in the table replacing $f_{j}^{\prime}(\phi_{j}^{k})$. All other entries in the table remain unchanged. The quantity $\phi_{j}^{k+1}$ is not explicitly stored.
\end{enumerate}
\end{minipage}}

Eliminating~$u^{k}$ recovers the update~\eqref{eq:w-update} for $x^{k}$. We now describe how the Finito \citep{finito} and MISO$\mu$ \citep{miso2} methods are closely related to SAGA. Both Finito and MISO$\mu$  use updates of the following form, for a step length $\stepsize$:
\vspace{-3mm}
\begin{equation}
	x^{k+1} = \frac{1}{n}\sum_{i} \phi^{k}_{i}  - \stepsize \sum^{n}_{i=1}f_{i}^{\prime}(\phi_{i}^{k}). \label{eq:finito}
\end{equation}
The step size used is of the order of $1/\mu n$. To simplify the discussion of this algorithm we will introduce the notation $\bar{\phi}=\frac{1}{n}\sum_{i} \phi^{k}_{i}$.

SAGA can be interpreted as Finito, but with the quantity $\bar{\phi}$ replaced with $u$, which is updated in the same way as $\bar{\phi}$, but \emph{in expectation}. To see this, consider how $\bar{\phi}$ changes in expectation:
\[
	\mathbb{E} \left[ \bar{\phi}^{k+1} \right]   =  \mathbb{E} \left[ \bar{\phi}^k + \frac{1}{n} \left( x^{k}  - \phi^{k}_{j}   \right) \right]
	 = \bar{\phi}^k + \frac{1}{n} \left( x^{k}  - \bar{\phi}^k   \right). 
\]
The update is identical in expectation to the update for $u$, $u^{k+1}=u^{k}+\frac{1}{n}(x^{k}-u^{k})$. 
There are three advantages of SAGA over Finito/MISO$\mu$. SAGA does not require strong convexity to work, it has support for proximal operators, and it does not require storing the $\phi_i$ values. MISO has proven support for proximal operators only in the case where impractically small step sizes are used \citep{miso2}.
The big advantage of Finito/MISO$\mu$ is that when using a per-pass re-permuted access ordering, empirical speed-ups of up-to a factor of 2x has been observed. This access order can also be used with the other methods discussed, but with smaller empirical speed-ups. Finito/MISO$\mu$ is particularly useful when $f_i$ is computationally expensive to compute compared to the extra storage costs required over the other methods.

\subsection*{SDCA}
\label{subsec:sdca}
The Stochastic Dual Coordinate Descent (SDCA) \citep{SDCA} method on the surface appears quite different from the other methods considered. It works with the convex conjugates of the $f_{i}$ functions. However, in this section we show a novel transformation of SDCA into an equivalent method that only works with primal quantities, and is closely related to the MISO$\mu$ method.

Consider the following algorithm:
\begin{framed}
\textbf{SDCA algorithm in the primal}

Step $k+1$:
\begin{enumerate}
  \setlength{\itemsep}{1pt}
\item Pick an index $j$ uniformly at random.
\item Compute $\phi_{j}^{k+1}=\text{prox}_{\stepsize}^{f_{j}}(z)$,
where $\stepsize=\frac{1}{\mu n}$ and
$z=-\stepsize  \sum^{n}_{i\neq j}f_{i}^{\prime}(\phi_{i}^{k})$.
\item Store the gradient $f_{j}^{\prime}(\phi_{j}^{k+1})=\frac{1}{\stepsize}\left(z-\phi_{j}^{k+1}\right)$
in the table at location $j$.
For $i\neq j$, the table entries are unchanged ($f_{i}^{\prime}(\phi_{i}^{k+1})=f_{i}^{\prime}(\phi_{i}^{k})$).
\end{enumerate}
At completion, return $x^{k}=-\stepsize \sum_{i}^{n}f_{i}^{\prime}(\phi_{i}^{k})$
 .
\end{framed}

We claim that this algorithm is equivalent to the version of SDCA where exact block-coordinate maximisation is used on the dual.\footnote{More precisely, to Option I of Prox-SDCA as described in~\cite[Figure 1]{sdca-accel}. We will simply refer to this method as ``SDCA'' in this paper for brevity.} Firstly, note that while SDCA was originally described for one-dimensional outputs (binary classification or regression), it has been expanded to cover the multi-class predictor case \citep{sdca-accel} (called Prox-SDCA there). In this case, the primal objective has a separate strongly convex regulariser, and the functions $f_i$ are restricted to the form $f_{i}(x) :=\psi_{i}(X_{i}^{T}x)$, where $X_{i}$ is a $d \times p$ feature matrix, and $\psi_{i}$ is the loss function that takes a $p$ dimensional input, for $p$ classes. To stay in the same general setting as the other incremental gradient methods, we work directly with the $f_i(x)$ functions rather than the more structured $\psi_{i}(X_{i}^{T}x)$. The dual objective to maximise then becomes 
\[ 
D(\alpha) = \left[
-\frac{\mu}{2} \left\Vert \frac{1}{\mu n} \sum_{i=1}^n \alpha_i \right\Vert^2 - \frac{1}{n} \sum_{i=1}^n f_i^*(-\alpha_i) \right],
\] where $\alpha_i$'s are $d$-dimensional dual variables. Generalising the exact block-coordinate maximisation update that SDCA performs to this form, we get the dual update for block $j$ (with $x^k$ the current primal iterate): 
\begin{equation} \label{eq:sdcaBlock}
\alpha_{j}^{k+1}=\alpha_{j}^{k}+\argmax_{\Delta a_{j} \in \mathbb{R}^d}\left\{ -f_{j}^{*}\left(-\alpha_{j}^{k}-\Delta\alpha_{j}\right)-\frac{\mu n}{2}\left\Vert x^{k}+\frac{1}{\mu n}\Delta\alpha_{j}\right\Vert ^{2} \right\}.
\end{equation}
In the special case where $f_{i}(x) =\psi_{i}(X_{i}^{T}x)$, we can see that~\eqref{eq:sdcaBlock} gives exactly the same update as Option~I of Prox-SDCA in~\cite[Figure 1]{sdca-accel}, which operates instead on the equivalent $p$-dimensional dual variables $\tilde{\alpha}_i$ with the relationship that $\alpha_i = X_i \tilde{\alpha}_i$.\footnote{This is because $f_i^*(\alpha_i) = \inf\limits_{\tilde{\alpha}_i \textrm{ s.t. } \alpha_i = X_i \tilde{\alpha}_i} \psi_i^*(\tilde{\alpha}_i)$.} As noted by Shalev-Shwartz \& Zhang  \citep{sdca-accel}, the update~\eqref{eq:sdcaBlock} is actually an instance of the proximal operator of the convex conjugate of $f_j$. Our primal formulation exploits this fact by using a relation between the proximal operator of a function and its convex conjugate known as the Moreau decomposition:
\[
\text{prox}^{f^{*}}(v)=v-\text{prox}^{f}(v).
\]
This decomposition allows us to compute the proximal operator of the conjugate via the primal proximal operator. As this is the only use in the basic SDCA method of the conjugate function, applying this decomposition allows us to completely eliminate the ``dual" aspect of the algorithm, yielding the above primal form of SDCA. The dual variables are related to the primal representatives $\phi_i$'s through $\alpha_i = -f^\prime_i(\phi_i)$. The KKT conditions ensure that if the $\alpha_i$ values are dual optimal then $x^k =  \stepsize \sum_i \alpha_i$ as defined above is primal optimal. The same trick is commonly used to interpret Dijkstra's set intersection as a primal algorithm instead of a dual block coordinate descent algorithm~\citep{prox-dk}.

The primal form of SDCA differs from the other incremental gradient methods described in this section in that it assumes strong convexity is induced by a separate strongly convex regulariser, rather than each $f_i$ being strongly convex. In fact, SDCA can be modified to work without a separate regulariser, giving a method that is at the midpoint between Finito and SDCA. We detail such a method in the supplementary material.

\subsection*{SDCA variants}
\label{subset:other-sdca}
%\vspace{-2mm}
The SDCA theory has been expanded to cover a number of other methods of performing the coordinate step \citep{sdca-accel}.  These variants replace the proximal operation in our primal interpretation in the previous section with an update where $\phi_{j}^{k+1}$ is chosen so that:
$
 	 f_{j}^{\prime} (\phi_{j}^{k+1})  = (1-\beta) f_{j}^{\prime} (\phi_{j}^{k}) + \beta f_{j}^{\prime} (x^{k})
$, 
where $x^{k}=-\frac{1}{\mu n}\sum_{i}f_{i}^{\prime}(\phi_{i}^{k})$. The variants differ in how $\beta \in [0,1]$  
is chosen. Note that $\phi_{j}^{k+1}$ does not actually have to be explicitly known, just the gradient $f_{j}^{\prime} (\phi_{j}^{k+1})$, which is the result of the above interpolation. Variant 5 by Shalev-Shwartz \& Zhang \citep{sdca-accel} does not require operations on the conjugate function, it simply uses $\beta = \frac{\mu n}{L+\mu n}$. The most practical variant performs a line search involving the convex conjugate to determine $\beta$. As far as we are aware, there is no simple primal equivalent of this line search. So in cases where we can not compute the proximal operator from the standard SDCA variant, we can either introduce a tuneable parameter into the algorithm ($\beta$), or use a dual line search, which requires an efficient way to evaluate the convex conjugates of each~$f_i$.

%\vspace{-2mm}
\section{Implementation}
\label{sec:impl}
\vspace{-2mm}
We briefly discuss some implementation concerns:
\begin{itemize}
%\vspace{-2mm}
\item For many problems each derivative $f^{\prime}_i$ is just a simple weighting of the $i$th data vector. Logistic regression and least squares have this property. In such cases, instead of storing the full derivative $f^{\prime}_i$ for each $i$, we need only to store the weighting constants. This reduces the storage requirements to be the same as the SDCA method in practice. A similar trick can be applied to multi-class classifiers with $p$ classes by storing $p-1$ values for each $i$.
\item Our algorithm assumes that initial gradients are known for each $f_i$ at the starting point $x^{0}$. Instead, a heuristic may be used where during the first pass, data-points are introduced one-by-one, in a non-randomized order, with averages computed in terms of those data-points processed so far. This procedure has been successfully used with SAG \citep{SAG}.
\item The SAGA update as stated is slower than necessary when derivatives are sparse. A just-in-time updating of $u$ or $x$ may be performed just as is suggested for SAG \citep{SAG}, which ensures that only sparse updates are done at each iteration.
\item We give the form of SAGA for the case where each $f_i$ is strongly convex. However in practice we usually have only convex $f_i$, with strong convexity in $f$ induced by the addition of a quadratic regulariser. This quadratic regulariser may be split amongst the $f_i$ functions evenly, to satisfy our assumptions. It is perhaps easier to use a variant of SAGA where the regulariser $\frac{\mu}{2}||x||^2$ is explicit, such as the following modification of Equation (\ref{eq:w-update}):
\[
x^{k+1}=\left(1-\stepsize \mu \right)x^{k}-\stepsize \left[ f_{j}^{\prime}(x^{k})-f_{j}^{\prime}(\phi_{j}^{k})+\frac{1}{n}\sum_{i}f_{i}^{\prime}(\phi_{i}^{k})\right].
 \]
 For sparse implementations instead of scaling $x^k$ at each step, a separate scaling constant $\beta^k$ may be scaled instead, with $\beta^k x^k$ being used in place of $x^k$. This is a standard trick used with stochastic gradient methods.
\end{itemize}
For sparse problems with a quadratic regulariser the just-in-time updating can be a little intricate. In the supplementary material we provide example python code showing a correct implementation that uses each of the above tricks.

\vspace{-2mm}
\section{Theory}
\label{sec:theory}
\vspace{-2mm}
In this section, all expectations
are taken with respect to the choice of $j$ at iteration $k+1$ and
conditioned on $x^{k}$ and each $f_{i}^{\prime}(\phi_{i}^{k})$
unless stated otherwise.

We start with two basic lemmas that just state properties of convex functions, followed by Lemma~\ref{lem:ip-bound}, which is specific to our algorithm. The proofs of each of these lemmas is in the supplementary material.

%%%%%%%%%%%%%%%%%%%%%%%%%%%%%%%%%%%%%%%%%

%%%%%%%%%%%%%%%%%%%%%%%%%%%%%%%%%%%%%%%%
\begin{lem}
\label{lem:ip-bound}
Let $f(x)=\frac{1}{n}\sum_{i=1}^{n}f_{i}(x)$. Suppose each $f_i$ is $\mu$-strongly convex and has Lipschitz continuous gradients with constant $L$. Then for all $x$ and $x^{*}$:
\begin{align*}
\left\langle f^{\prime}(x),x^{*}-x\right\rangle  \leq {} &
\frac{L-\mu}{L}\left[f(x^{*})-f(x)\right]-\frac{\mu}{2}\left\Vert x^{*}-x\right\Vert ^{2} \\
& -\frac{1}{2Ln}\sum_{i}\left\Vert f_{i}^{\prime}(x^{*})-f_{i}^{\prime}(x)\right\Vert ^{2}-\frac{\mu}{L}\left\langle f^{\prime}(x^{*}),x-x^{*}\right\rangle. 
\end{align*}
\end{lem}
%%%%%%%%%%%%%%%%%%%%%%%%%%%%%%%%%%%%%%%%%%
\begin{lem}
\label{lem:grad-diff-phii}We have that for all $\phi_{i}$ and $x^{*}$:
\[
\frac{1}{n}\sum_{i}\left\Vert f_{i}^{\prime}(\phi_{i})-f_{i}^{\prime}(x^{*})\right\Vert ^{2}
\leq 2L\left[\frac{1}{n}\sum_{i}f_{i}(\phi_{i})-f(x^{*})-\frac{1}{n}\sum_{i}\left\langle f_{i}^{\prime}(x^{*}),\phi_{i}-x^{*}\right\rangle \right].
\]
\end{lem}

%%%%%%%%%%%%%%%%%%%%%%%%%%%%%%%%%%%%%%%%%%%
\begin{lem}
\label{lem:wchange} It holds that for any $\phi_{i}^{k}$, $x^{*}$, $x^k$
and $\beta>0$, with $w^{k+1}$ as defined in Equation \ref{eq:main-update}:
\begingroup\makeatletter\def\f@size{9}\check@mathfonts
\begin{align*}
\mathbb{E}\left\Vert w^{k+1}-x^k-\stepsize f^{\prime}(x^{*})\right\Vert ^{2} 
\leq {} & \stepsize^2 (1+\beta^{-1}) \mathbb{E}\left\Vert f_{j}^{\prime}(\phi_{j}^{k})-f_{j}^{\prime}(x^{*})\right\Vert ^{2}
+ \stepsize^2 (1+\beta) \mathbb{E}\left\Vert f_{j}^{\prime}(x^k)-f_{j}^{\prime}(x^{*})\right\Vert ^{2} \\
&- \stepsize^2 \beta \left\Vert f^{\prime}(x^k)-f^{\prime}(x^{*}) \right\Vert ^{2}.
\end{align*}
\endgroup
\end{lem}

%%%%%%%%%%%%%%%%%%%%%%%%%%%%%%%%%%%%%%%%%%%%%%%
\begin{thm}
\label{thm:strong-convex-thm} With $x^*$ the optimal solution, define the Lyapunov function $T$ as: 
\[
T^k := T(x^k, \{\phi_i^k\}_{i=1}^n) :=\frac{1}{n}\sum_{i}f_{i}(\phi^k_{i})-f(x^{*})-\frac{1}{n}\sum_{i}\left\langle f_{i}^{\prime}(x^{*}),\phi^k_{i}-x^{*}\right\rangle +c\left\Vert x^k-x^{*}\right\Vert ^{2}.
\]
 Then with $\stepsize=\frac{1}{2(\mu n+L)}$, $c=\frac{1}{2\stepsize(1-\stepsize\mu)n}$, and $\kappa=\frac{1}{\stepsize\mu}$, we have the following expected change in the Lyapunov function between steps of the SAGA algorithm (conditional on $T^k$):
\[
\mathbb{E}[T^{k+1}]\leq(1-\frac{1}{\kappa})T^{k}.
\]
\end{thm}
\begin{proof}
\vspace{-2mm}
The first three terms in $T^{k+1}$ are straight-forward to simplify:
\begin{align*}
\mathbb{E}\left[\frac{1}{n}\sum_{i}f_{i}(\phi_{i}^{k+1})\right] & = \frac{1}{n}f(x^k)+\left(1-\frac{1}{n}\right)\frac{1}{n}\sum_{i}f_{i}(\phi_{i}^{k}). \\
\mathbb{E}\left[-\frac{1}{n}\sum_{i}\left\langle f_{i}^{\prime}(x^{*}),\phi_{i}^{k+1}-x^{*}\right\rangle \right] & = 
 -\frac{1}{n}\left\langle f^{\prime}(x^{*}),x^k-x^{*}\right\rangle \!- \! \left(1 \! - \! \frac{1}{n}\right)\frac{1}{n}\sum_{i}\left\langle f_{i}^{\prime}(x^{*}),\phi_{i}^{k}-x^{*}\right\rangle .
\end{align*}
For the change in the last term of $T^{k+1}$, we apply the non-expansiveness of the proximal operator\footnote{Note that the first equality below is the only place in the proof where we use the fact that $x^*$ is an optimality point.}:
\begin{align*}
c\left\Vert x^{k+1}-x^{*}\right\Vert ^{2} & = c\left\Vert \text{prox}_{\stepsize}(w^{k+1})-\text{prox}_{\stepsize}(x^{*}-\stepsize f^{\prime}(x^{*}))\right\Vert ^{2}\\
 & \leq c\left\Vert w^{k+1}-x^{*}+\stepsize f^{\prime}(x^{*})\right\Vert ^{2}.
\end{align*}
We expand the quadratic and apply $\mathbb{E}[w^{k+1}]=x^{k}-\stepsize f^{\prime}(x^{k})$
to simplify the inner product term:
\begin{align*}
  & \!\!\!\!\! c\mathbb{E}\left\Vert w^{k+1}-x^{*}+\stepsize f^{\prime}(x^{*})\right\Vert ^{2} \, = \,  c\mathbb{E}\left\Vert x^k-x^{*}+w^{k+1}-x^k+\stepsize f^{\prime}(x^{*})\right\Vert ^{2}\\
 = {} &  c\left\Vert x^k-x^{*}\right\Vert ^{2}+2c\mathbb{E}\left[\left\langle w^{k+1}-x^k+\stepsize f^{\prime}(x^{*}),x^k-x^{*}\right\rangle \right]+c\mathbb{E}\left\Vert w^{k+1}-x^k+\stepsize f^{\prime}(x^{*})\right\Vert ^{2}\\
 = {} &  c\left\Vert x^k-x^{*}\right\Vert ^{2}-2c\stepsize\left\langle f^{\prime}(x^k)-f^{\prime}(x^{*}),x^k-x^{*}\right\rangle +c\mathbb{E}\left\Vert w^{k+1}-x^k+\stepsize f^{\prime}(x^{*})\right\Vert ^{2}\\
 \leq {} &  c\left\Vert x^k-x^{*}\right\Vert ^{2}-2c\stepsize\left\langle f^{\prime}(x^k),x^k-x^{*}\right\rangle +2c \stepsize \left\langle f^{\prime}(x^{*}),x^k-x^{*}\right\rangle 
 - c \stepsize^{2} \beta \left\Vert f^{\prime}(x^k)-f^{\prime}(x^{*}) \right\Vert ^{2}\\
 &   +\left(1+\beta^{-1}\right)c \stepsize^{2}\mathbb{E}\left\Vert f_{j}^{\prime}(\phi_{j}^{k})-f_{j}^{\prime}(x^{*})\right\Vert ^{2}+\left(1+\beta\right)c \stepsize^{2}\mathbb{E}\left\Vert f_{j}^{\prime}(x^k)-f_{j}^{\prime}(x^{*})\right\Vert ^{2}.\quad\text{(Lemma \ref{lem:wchange})}
\end{align*}
The value of $\beta$ shall be fixed later. Now we apply Lemma \ref{lem:ip-bound} to bound $-2c\stepsize\left\langle f^{\prime}(x^k),x^k-x^{*}\right\rangle $ and Lemma \ref{lem:grad-diff-phii} to bound $\mathbb{E}\left\Vert f_{j}^{\prime}(\phi_{j}^{k})-f_{j}^{\prime}(x^{*})\right\Vert ^{2}$:
\begin{align*}
c\mathbb{E}\left\Vert x^{k+1}-x^{*}\right\Vert ^{2} &\leq \left(c-c\stepsize\mu \right)\left\Vert x^k-x^{*}\right\Vert ^{2}+\left((1+\beta)c\stepsize^{2}-\frac{c\stepsize }{L}\right)\mathbb{E}\left\Vert f_{j}^{\prime}(x^k)-f_{j}^{\prime}(x^{*})\right\Vert ^{2}\\
  & -\frac{2c\stepsize(L-\mu)}{L}\left[f(x^k)-f(x^{*})-\left\langle f^{\prime}(x^{*}),x^k-x^{*}\right\rangle \right]
      - c\stepsize^{2}\beta\left\Vert f^{\prime}(x^k)-f^{\prime}(x^{*}) \right\Vert ^{2} \\
 &   +2\left(1+\beta^{-1}\right)c\stepsize^2 L\left[\frac{1}{n}\sum_{i}f_{i}(\phi_{i}^{k})-f(x^{*})-\frac{1}{n}\sum_{i}\left\langle f_{i}^{\prime}(x^{*}),\phi_{i}^{k}-x^{*}\right\rangle \right].
\end{align*}
We can now combine the bounds that we have derived for each term in $T$,
and pull out a fraction $\frac{1}{\kappa}$ of $T^{k}$ (for any $\kappa$ at this point). Together with the inequality 
$-\left\Vert f^{\prime}(x^k)-f^{\prime}(x^{*})\right\Vert ^{2}\leq-2\mu\left[f(x^k)-f(x^{*})-\left\langle f^{\prime}(x^{*}),x^k-x^{*}\right\rangle \right]$~\cite[Thm. 2.1.10]{nes-book}, that yields:
\begingroup\makeatletter\def\f@size{9}\check@mathfonts
\begin{align}
\mathbb{E}[T^{k+1}]&-T^{k} \leq -\frac{1}{\kappa}T^{k} + \left(\frac{1}{n}-\frac{2c\stepsize(L-\mu)}{L} -2c\stepsize^2\mu\beta \right)\left[f(x^k)-f(x^{*})-\left\langle f^{\prime}(x^{*}),x^k-x^{*}\right\rangle \right]\nonumber \\
 &   +\left(\frac{1}{\kappa}+2(1+\beta^{-1})c\stepsize^{2}L-\frac{1}{n}\right)\left[\frac{1}{n}\sum_{i}f_{i}(\phi_{i}^{k})-f(x^{*})-\frac{1}{n}\sum_{i}\left\langle f_{i}^{\prime}(x^{*}),\phi_{i}^{k}-x^{*}\right\rangle \right]\nonumber \\
 &   +\left(\frac{1}{\kappa}-\stepsize\mu\right)c\left\Vert x^k-x^{*}\right\Vert ^{2}+\left((1+\beta)\stepsize-\frac{1}{ L}\right) c\stepsize \mathbb{E}\left\Vert f_{j}^{\prime}(x^k)-f_{j}^{\prime}(x^{*})\right\Vert ^{2} . \label{eq:constants-strong}
\end{align}
\endgroup
Note that each of the terms in square brackets are positive, and it
can be readily verified that our assumed values for the constants
($\stepsize=\frac{1}{2(\mu n+L)}$, $c=\frac{1}{2\stepsize(1-\stepsize\mu)n}$, and $\kappa=\frac{1}{\stepsize\mu}$), together with $\beta=\frac{2\mu n+L}{L}$ ensure
that each of the quantities in round brackets are non-positive (the constants were determined by setting all the round brackets to zero except the second one  --- see~\citep{adefazio-thesis2014} for the details). \\
\textbf{Adaptivity to strong convexity result:} Note that when using the $\stepsize=\frac{1}{3L}$ step size, the same $c$ as above can be used with $\beta=2$ and $\frac{1}{\kappa}=\min\left\{ \frac{1}{4n},\,\frac{\mu}{3L}\right\}$ to ensure non-positive terms.
\end{proof}
%%%%%%%%%%%%%%%%%%%%%%%%%%%%%%%%
\begin{cor}
Note that $c\left\Vert x^{k}-x^{*}\right\Vert ^{2}\leq T^{k}$, and therefore
by chaining the expectations, plugging in the constants explicitly and using $\mu (n-0.5) \leq \mu n$ to simplify the expression, we get:
\begingroup\makeatletter\def\f@size{9}\check@mathfonts
\[
\mathbb{E}\left[\left\Vert x^{k}-x^{*}\right\Vert ^{2}\right]\leq\left(1-\frac{\mu}{2(\mu n+L)}\right)^{k}\left[\left\Vert x^{0}-x^{*}\right\Vert ^{2} 
+ \frac{n}{\mu n + L}
\left[f(x^{0}) -\left\langle f^{\prime}(x^{*}),x^{0}-x^{*}\right\rangle  -f(x^{*})\right]\right].
\]
\endgroup
Here the expectation is over all choices of index $j^{k}$ up to step $k$. 
\end{cor}

\begin{figure}[t]
\begin{minipage}{5mm}
\begin{sideways}\small Function sub-optimality \end{sideways}
\end{minipage}
\begin{centering}
\begin{minipage}{135mm} \centering
\includegraphics[scale=0.96]{mnist-stats-runOn-2014-06-05--17-44-30.pdf} 
\includegraphics[scale=0.96]{covtype-stats-runOn-2014-06-05--17-52-28.pdf}
\includegraphics[scale=0.96]{ijcnn1-stats-runOn-2014-06-05--17-53-16.pdf}
\includegraphics[scale=0.96]{millionsong-stats-runOn-2014-06-05--19-28-10.pdf}
\includegraphics[scale=0.96]{mnist-nl1-stats-runOn-2014-06-05--23-55-35.pdf} 
\includegraphics[scale=0.96]{covtype-nl1-stats-runOn-2014-06-06--11-28-36.pdf}
\includegraphics[scale=0.96]{ijcnn1-nl1-stats-runOn-2014-06-05--19-43-20.pdf}
\includegraphics[scale=0.96]{millionsong-nl1-stats-runOn-2014-06-06--11-22-14.pdf}
\end{minipage}
\end{centering}
\center
{\small Gradient evaluations / $n$}

\includegraphics[scale=1.5,trim=0 0 30 0, clip]{caption.pdf}
\vskip -0.1in
\caption{\small From left to right we have the MNIST, COVTYPE, IJCNN1 and MILLIONSONG datasets. Top row is the L2 regularised case, bottom row the  L1 regularised case. \label{fig:plots}}
\vskip -0.2in
\end{figure}
\vspace{-2mm}
\section{Experiments}
\label{sec:experiments}
\vspace{-3mm}
We performed a series of experiments to validate the effectiveness of SAGA. We tested a binary classifier on MNIST, COVTYPE, IJCNN1 and a least squares predictor on MILLIONSONG. Details of these datasets can be found in \citep{finito}. We used the same code base for each method, just changing the main update rule. SVRG was tested with the recalibration pass used every $n$ iterations, as suggested in \citep{semi}. Each method had its step size parameter chosen so as to give the fastest convergence.

We tested with a L2 regulariser, which all methods support, and with a L1 regulariser on a subset of the methods. The results are shown in Figure \ref{fig:plots}. We can see that Finito (perm) performs the best on a per epoch equivalent basis, but it can be the most expensive method per step. SVRG is similarly fast on a per epoch basis, but when considering the number of gradient evaluations per epoch is double that of the other methods for this problem, it is middle of the pack. SAGA can be seen to perform similar to the non-permuted Finito case, and to SDCA. Note that SAG is slower than the other methods at the beginning. To get the optimal results for SAG, an adaptive step size rule needs to be used rather than the constant step size we used. In general, these tests confirm that the choice of methods should be done based on their properties as discussed in Section \ref{sec:related-work}, rather than their convergence rate. 

\nocite{prox-svrg}
\nocite{nes-book}

\pagebreak
\bibliography{nips2014}
\bibliographystyle{unsrt}


\newpage

\appendix

\section{The SDCA/Finito Midpoint Algorithm}
Using Lagrangian duality theory, SDCA can be shown at step $k$ as minimising the following lower bound:
\[
A^{k}(x)=\frac{1}{n}f_{j}(x) +
 \frac{1}{n}\sum^{n}_{i \neq j}\left[f_{i}(\phi_{i}^{k})+\left\langle f_{i}^{\prime}(\phi_{i}^{k}),x-\phi_{i}^{k}\right\rangle \right]+\frac{\mu}{2}\left\Vert x\right\Vert ^{2}.
\]
Instead of directly including the regulariser in this bound, we can use the standard strong convexity lower bound for each $f_i$, by removing $\frac{\mu}{2}\left\Vert x\right\Vert ^{2}$ and changing the expression in the summation to 
$f_{i}(\phi_{i}^{k})+\left\langle f_{i}^{\prime}(\phi_{i}^{k}),x-\phi_{i}^{k}\right\rangle +\frac{\mu}{2}\left\Vert x-\phi_{i}\right\Vert ^{2}$.
The transformation to having strong convexity within the $f_{i}$ functions yields the following simple modification to the algorithm:
 $\phi_{j}^{k+1}=\text{prox}_{(\mu(n-1))^{-1}}^{f_{j}}(z)$, where:
\[
z=\frac{1}{n-1}\sum_{i\neq j}\phi_{i}^{k}-\frac{1}{\mu(n-1)}\sum_{i\neq j}f_{i}^{\prime}(\phi_{i}^{k}).
\]
It can be shown that after this update:
\[
	x^{k+1}=\phi_{j}^{k+1}=\frac{1}{n}\sum_{i}\phi_{i}^{k+1}-\frac{1}{\mu n}\sum_{i}f_{i}^{\prime}(\phi_{i}^{k+1}).
\]
Now the similarity to Finito is apparent if this equation is compared Equation~8: %\ref{eq:finito}: -- SLJ: need to keep track of this when changing!
$x^{k+1} = \frac{1}{n}\sum_{i} \phi^{k}_{i}  - \stepsize \sum^{n}_{i=1}f_{i}^{\prime}(\phi_{i}^{k})$. The only difference is that the vectors on the right hand side of the equation are at their values at step $k+1$ instead of $k$. Note that there is a circular dependency here, as $\phi_{j}^{k+1}:=x^{k+1}$ but $\phi_{j}^{k+1}$ appears in the definition of $x^{k+1}$. Solving the proximal operator is the resolution of the circular dependency. This mid-point between Finito and SDCA is interesting in it's own right, as it appears experimentally to have similar robustness to permuted orderings as Finito, but it has no tunable parameters like SDCA.

When the proximal operator above is fast to compute, say on the same order as just evaluating $f_j$, then SDCA can be the best method among those discussed. It is a little slower than the other methods discussed here, but it has no tunable parameters at all. It is also the only choice when each $f_i$ is not differentiable. The major disadvantage of SDCA is that it can not handle non-strongly convex problems directly. Although like most methods, adding a small amount of quadratic regularisation can be used to recover a convergence rate. It is also not adapted to use proximal operators \emph{for the regulariser} in the composite objective case. The requirement of computing the proximal operator of each loss $f_i$ initially appears to be a big disadvantage, however there are variants of SDCA that remove this requirement, but they introduce additional downsides.

\section{Lemmas}

%appendixlem is numbered A1, so it doesn't clash with the main document numbering.
\begin{lem}
\label{thm:strong-lb}Let $f$ be $\mu$-strongly convex and have Lipschitz continuous gradients with constant $L$. Then we have for all $x$ and $y$:
\begin{align*}
f(x) \geq {} & f(y)+\left\langle f^{\prime}(y),x-y\right\rangle +\frac{1}{2\left(L-\mu \right)}\left\Vert f^{\prime}(x)-f^{\prime}(y)\right\Vert ^{2} \\
& +\frac{\mu L}{2\left(L-\mu \right)}\left\Vert y-x\right\Vert ^{2}+\frac{\mu}{\left(L-\mu \right)}\left\langle f^{\prime}(x)-f^{\prime}(y),y-x\right\rangle .
\end{align*}
\end{lem}
\begin{proof}
Define the function $g$ as $g(x)=f(x)-\frac{\mu}{2}\left\Vert x\right\Vert ^{2}$.
Then the gradient is $g^{\prime}(x)=f^{\prime}(x)-\mu x$. $g$~has a
Lipschitz gradient with constant $L-\mu$. By convexity, we have~\cite[Thm. 2.1.5]{nes-book}:
\[
g(x)\geq g(y)+\left\langle g^{\prime}(y),x-y\right\rangle +\frac{1}{2(L-\mu)}\left\Vert g^{\prime}(x)-g^{\prime}(y)\right\Vert ^{2}.
\]
Substituting in the definition of $g$ and $g^{\prime},$ and simplifying the
terms gives the result.\end{proof}

\begin{lem}
Let $f(x)=\frac{1}{n}\sum_{i=1}^{n}f_{i}(x)$. Suppose each $f_i$ is $\mu$-strongly convex and has Lipschitz continuous gradients with constant $L$. Then for all $x$ and $x^{*}$:
\begingroup\makeatletter\def\f@size{9}\check@mathfonts
\[
\left\langle f^{\prime}(x),x^{*}-x\right\rangle  \leq 
\frac{L-\mu}{L}\left[f(x^{*})-f(x)\right]-\frac{\mu}{2}\left\Vert x^{*}-x\right\Vert ^{2}-\frac{1}{2Ln}\sum_{i}\left\Vert f_{i}^{\prime}(x^{*})-f_{i}^{\prime}(x)\right\Vert ^{2}-\frac{\mu}{L}\left\langle f^{\prime}(x^{*}),x-x^{*}\right\rangle. 
\]
\endgroup
\end{lem}
\begin{proof}
This is a straight-forward corollary of Lemma \ref{thm:strong-lb}, using $y=x^*$, and averaging over the $f_i$ functions.
\end{proof}

\begin{lem}
\label{lem:grad-diff-phii}
We have that for all $\phi_{i}$ and $x^{*}$:
\[
\frac{1}{n}\sum_{i}\left\Vert f_{i}^{\prime}(\phi_{i})-f_{i}^{\prime}(x^{*})\right\Vert ^{2}
\leq 2L\left[\frac{1}{n}\sum_{i}f_{i}(\phi_{i})-f(x^{*})-\frac{1}{n}\sum_{i}\left\langle f_{i}^{\prime}(x^{*}),\phi_{i}-x^{*}\right\rangle \right].
\]
\end{lem}
\begin{proof}
Apply the standard inequality $f(y)\geq f(x)+\left\langle f^{\prime}(x),y-x\right\rangle +\frac{1}{2L}\left\Vert f^{\prime}(x)-f^{\prime}(y)\right\Vert ^{2}$,
with $y=\phi_{i}$ and $x=x^{*}$, for each $f_{i}$, and sum. \end{proof}

\begin{lem}
\label{lem:wchange}
It holds that for any $\phi_{i}^{k}$, $x^{*}$, $x^k$
and $\beta>0$, with $w^{k+1}$ as defined in Equation~1: %\ref{eq:main-update}:
\begin{align*}
\mathbb{E}\left\Vert w^{k+1}-x^k-\stepsize f^{\prime}(x^{*})\right\Vert ^{2} 
\leq {} & \stepsize^2 (1+\beta^{-1}) \mathbb{E}\left\Vert f_{j}^{\prime}(\phi_{j}^{k})-f_{j}^{\prime}(x^{*})\right\Vert ^{2}
+ \stepsize^2 (1+\beta) \mathbb{E}\left\Vert f_{j}^{\prime}(x^k)-f_{j}^{\prime}(x^{*})\right\Vert ^{2} \\
&- \stepsize^2 \beta \left\Vert f^{\prime}(x^k)-f^{\prime}(x^{*}) \right\Vert ^{2}.
\end{align*}
\end{lem}
\begin{proof}
We follow a similar argument as occurs in the SVRG proof \citep{svrg} for this term, but with a tighter argument. The tightening comes from using $\left\Vert x+y\right\Vert ^{2}\leq(1+\beta^{-1})\left\Vert x\right\Vert ^{2}+(1+\beta)\left\Vert y\right\Vert ^{2}$ instead of the simpler  $\beta=1$ case they use.
The other key trick is the use of the standard variance decomposition $\mathbb{E}[\left\Vert X-\mathbb{E}[X]\right\Vert ^{2}]=\mathbb{E}[\left\Vert X\right\Vert ^{2}]-\left\Vert \mathbb{E}[X]\right\Vert ^{2}$ three times.
\begingroup\makeatletter\def\f@size{9}\check@mathfonts
\begin{align*}
 &  \mathbb{E}\left\Vert w^{k+1}-x^k+\stepsize f^{\prime}(x^{*})\right\Vert ^{2} \\
  =  {} & \mathbb{E}\bigg\Vert \underbrace{-\frac{\stepsize}{ n}\sum_{i}f_{i}^{\prime}(\phi_{i}^{k})+\stepsize  f^{\prime}(x^{*})+\stepsize  \left[f_{j}^{\prime}(\phi_{j}^{k})-f_{j}^{\prime}(x^k)\right]}_{:= \,\,\stepsize X}\bigg\Vert ^{2}\\
  =  {} & \stepsize^{2} \mathbb{E}\Bigg\Vert \overbrace{\Bigg[f_{j}^{\prime}(\phi_{j}^{k})-f_{j}^{\prime}(x^{*})-\frac{1}{n}\sum_{i}f_{i}^{\prime}(\phi_{i}^{k})+f^{\prime}(x^{*})\Bigg]
 -\Bigg[f_{j}^{\prime}(x^k)-f_{j}^{\prime}(x^{*}) }^{X} - \overbrace{ f^{\prime}(x^k) + f^{\prime}(x^{*})\Bigg]}^{\mathbb{E}[X]} \Bigg\Vert ^{2}  
 + \stepsize^{2} \bigg\Vert \overbrace{f^{\prime}(x^k)-f^{\prime}(x^{*})}^{\mathbb{E}[X]} \bigg\Vert ^{2} \\
 \leq {} & \stepsize^{2}(1+\beta^{-1})\mathbb{E}\left\Vert f_{j}^{\prime}(\phi_{j}^{k})-f_{j}^{\prime}(x^{*})-\frac{1}{n}\sum_{i}f_{i}^{\prime}(\phi_{i}^{k})+f^{\prime}(x^{*})\right\Vert ^{2} \\
 & + \stepsize^{2}(1+\beta) \mathbb{E}\left\Vert f_{j}^{\prime}(x^k)-f_{j}^{\prime}(x^{*}) 
       - f^{\prime}(x^k) + f^{\prime}(x^{*}) \right\Vert ^{2}
    + \stepsize^{2}\left\Vert f^{\prime}(x^k)-f^{\prime}(x^{*}) \right\Vert ^{2}\\
 & \textrm{\qquad \qquad (use variance decomposition twice more):} \\
 \leq {} & \stepsize^{2}(1+\beta^{-1})\mathbb{E}\left\Vert f_{j}^{\prime}(\phi_{j}^{k})-f_{j}^{\prime}(x^{*})\right\Vert ^{2}+\stepsize^{2}(1+\beta)\mathbb{E}\left\Vert f_{j}^{\prime}(x^k)-f_{j}^{\prime}(x^{*})\right\Vert ^{2} - \stepsize^{2} \beta \left\Vert f^{\prime}(x^k)-f^{\prime}(x^{*}) \right\Vert ^{2}.
\end{align*}
\endgroup
\end{proof}



\section{Non-strongly-convex Problems}

\begin{thm}
When each $f_{i}$ is convex, using $\stepsize=\frac{1}{3L}$,
we have for $\bar{x}^{k}=\frac{1}{k}\sum_{t=1}^{k}x^{t}$ that:
\[
\mathbb{E}\left[F(\bar{x}^{k})\right]-F(x^{*})\leq\frac{4n}{k}\left[\frac{2L}{n}\left\Vert x^{0}-x^{*}\right\Vert ^{2}+f(x^{0})-\left\langle f^{\prime}(x^{*}),x^{0}-x^{*}\right\rangle -f(x^{*})\right].
\]
Here the expectation is over all choices of index $j^{k}$ up to step $k$. 
\end{thm}
\begin{proof}
A more detailed version of this proof is available in \citep{adefazio-thesis2014}.
We proceed by using a similar argument as in Theorem~1, %\ref{thm:strong-convex-thm},
but we add an additional $\alpha\left\Vert x^k-x^{*}\right\Vert ^{2}$
together with the existing $c\left\Vert x^k-x^{*}\right\Vert ^{2}$
term in the Lyapunov function.

We will bound $\alpha\left\Vert x^k-x^{*}\right\Vert ^{2}$ in a different
manner to $c\left\Vert x^k-x^{*}\right\Vert ^{2}$. Define $\Delta=-\frac{1}{\stepsize}\left(w^{k+1}-x^k\right)-f^{\prime}(x^k)$,
the difference between our approximation to the gradient at $x^k$ and
true gradient. Then instead of using the non-expansiveness property
at the beginning, we use a result proved for prox-SVRG~\cite[2nd eq. on p.12]{prox-svrg}:
\[
\alpha \mathbb{E}\left\Vert x^{k+1}-x^{*}\right\Vert ^{2}\leq\alpha\left\Vert x^k-x^{*}\right\Vert ^{2}-2\alpha\stepsize\mathbb{E}\left[F(x^{k+1})-F(x^{*})\right]+2\alpha\stepsize^{2}\mathbb{E}\left\Vert \Delta\right\Vert ^{2}.
\]
Although their quantity $\Delta$ is different, they only use the
property that $\mathbb{E}[\Delta]=0$ to prove the above equation. A full proof of this property for the SAGA algorithm that follows their argument appears in \citep{adefazio-thesis2014}.

To bound the $\Delta$ term, a small modification of the argument in Lemma \ref{lem:wchange} can be used, giving:
\[
\mathbb{E}\left\Vert \Delta\right\Vert ^{2} \leq \left(1+\beta^{-1}\right)\mathbb{E}\left\Vert f_{j}^{\prime}(\phi_{j}^{k})-f_{j}^{\prime}(x^{*})\right\Vert ^{2}+\left(1+\beta\right)\mathbb{E}\left\Vert f_{j}^{\prime}(x^{k})-f_{j}^{\prime}(x^{*})\right\Vert ^{2}.
\]
 Applying this gives:
\begin{align*}
\alpha \mathbb{E}\left\Vert x^{k+1}-x^{*}\right\Vert ^{2} \leq {} & \alpha\left\Vert x^k-x^{*}\right\Vert ^{2}-2\alpha\stepsize\mathbb{E}\left[F(x^{k+1})-F(x^{*})\right]\\
& +2(1+\beta^{-1})\alpha\stepsize^2\mathbb{E}\left\Vert f_{j}^{\prime}(\phi_{j}^{k})-f_{j}^{\prime}(x^{*})\right\Vert ^{2}+2\left(1+\beta\right)\alpha\stepsize^{2}\mathbb{E}\left\Vert f_{j}^{\prime}(x^k)-f_{j}^{\prime}(x^{*})\right\Vert ^{2}.
\end{align*}
As in Theorem~1, %\ref{thm:strong-convex-thm}, 
we then apply Lemma \ref{lem:grad-diff-phii} to bound $\mathbb{E}\left\Vert f_{j}^{\prime}(\phi_{j}^{k})-f_{j}^{\prime}(x^{*})\right\Vert ^{2}$. Combining with the rest of the Lyapunov function as was derived in Theorem~1 gives (we basically add the $\alpha$ terms to inequality~(10) with $\mu=0$):
\begin{eqnarray*}
 &  & \mathbb{E}[T^{k+1}]-T^{k} \\
 & \leq & \left(\frac{1}{n}-2c\gamma\right)\left[f(x^{k})-f(x^{*})-\left\langle f^{\prime}(x^{*}),x^{k}-x^{*}\right\rangle \right]-2\alpha\gamma\mathbb{E}\left[F(x^{k+1})-F(x^{*})\right]\\
 &  & +\left(4(1+\beta^{-1})\alpha L\gamma^{2}+2(1+\beta^{-1})cL\gamma^{2}-\frac{1}{n}\right)\left[\frac{1}{n}\sum_{i}f_{i}(\phi_{i}^{k})-f(x^{*})-\frac{1}{n}\sum_{i}\left\langle f_{i}^{\prime}(x^{*}),\phi_{i}^{k}-x^{*}\right\rangle \right] \\
 &  & +\left((1+\beta)c\gamma+2(1+\beta)\alpha\gamma-\frac{c}{L}\right)\gamma\mathbb{E}\left\Vert f_{j}^{\prime}(x^{k})-f_{j}^{\prime}(x^{*})\right\Vert ^{2}.
\end{eqnarray*}
As before, the terms in square brackets are positive by convexity. Given that our choice of step size is $\gamma=\frac{1}{3L}$ (to match the adaptive to strong convexity step size), we can set the three round brackets to zero by using $\beta=1$, $c=\frac{3L}{2n}$ and $\alpha=\frac{3L}{8n}$. We thus obtain:
\[
\mathbb{E}[T^{k+1}]-T^{k}\leq-\frac{1}{4n}\mathbb{E}\left[F(x^{k+1})-F(x^{*})\right].
\]
These expectations are conditional on information from step $k$. We now take the expectation with respect to all previous steps, yielding $\mathbb{E}[T^{k+1}]-\mathbb{E}[T^{k}] \leq-\frac{1}{4n}\mathbb{E}\left[F(x^{k+1})-F(x^{*})\right]$, where all expectations are unconditional. Further negating and summing for $k$ from $0$ to $k-1$ results in telescoping of the $T$ terms, giving:
\[
\frac{1}{4n}\mathbb{E}\left[\sum_{t=1}^{k}\left[F(x^{t})-F(x^{*})\right]\right]  \leq  T^{0}-\mathbb{E}[T^{k}].
\]
We can drop the $-\mathbb{E}\left[T^{k}\right]$ term since $T^{k}$ is
always positive. Then we apply convexity to pull the summation inside
of $F$, and multiply through by $4n/k$, giving:
\[
\mathbb{E}\left[F(\frac{1}{k}\sum_{t=1}^{k}x^{t})-F(x^{*})\right]
\leq \frac{1}{k}\mathbb{E}\left[\sum_{t=1}^{k}\left[F(x^{t})-F(x^{*})\right]\right]
\leq \frac{4n}{k}T^{0}.
\]
We get a $(c+\alpha) = \frac{15L}{8n} \leq \frac{2L}{n}$ term that we use in $T^{0}$ for simplicity.
\end{proof}

\section{Example Code for Sparse Least Squares \& Ridge Regression}
The SAGA method is quite easy to implement for dense gradients, however the implementation for sparse gradient problems can be tricky.
The main complication is the need for just-in-time updating of the elements of the iterate vector. This is needed to avoid having
to do any full dense vector operations at each iteration.
We provide below a simple implementation for the case of least-squares problems that illustrates how to correctly do this. The code is in the compiled Python (Cython) language.

\begin{python}
import random
import numpy as np
cimport numpy as np

cimport cython
from cython.view cimport array as cvarray

# Performs the lagged update of x by g.
cdef inline lagged_update(long k, double[:] x, double[:] g, unsigned long[:] lag, 
                          long[:] yindices, int ylen, double[:] lag_scaling, double a):
    
    cdef unsigned int i
    cdef long ind
    cdef unsigned long lagged_amount = 0
    
    for i in range(ylen):
        ind = yindices[i]
        lagged_amount = k-lag[ind]
        lag[ind] = k
        x[ind] += lag_scaling[lagged_amount]*(a*g[ind])
        

# Performs x += a*y, where x is dense and y is sparse.
cdef inline add_weighted(double[:] x, double[:] ydata , long[:] yindices, int ylen, double a):
    cdef unsigned int i
    
    for i in range(ylen):
        x[yindices[i]] += a*ydata[i]

# Dot product of a dense vector with a sparse vector
cdef inline spdot(double[:] x, double[:] ydata , long[:] yindices, int ylen):
    cdef unsigned int i
    cdef double v = 0.0
    
    for i in range(ylen):
        v += ydata[i]*x[yindices[i]]
        
    return v

def saga_lstsq(A,  double[:] b, unsigned int maxiter, props):
    
    # temporaries
    cdef double[:] ydata
    cdef long[:] yindices
    cdef unsigned int i, j, epoch, lagged_amount
    cdef long indstart, indend, ylen, ind
    cdef double cnew, Aix, cchange, gscaling
    
    # Data points are stored in columns in CSC format.
    cdef double[:] data = A.data
    cdef long[:] indices = A.indices
    cdef long[:] indptr = A.indptr
    
    cdef unsigned int m = A.shape[0] # dimensions
    cdef unsigned int n = A.shape[1] # datapoints
    
    cdef double[:] xk = np.zeros(m)
    cdef double[:] gk = np.zeros(m)
    
    cdef double eta = props['eta'] # Inverse step size = 1/gamma
    cdef double reg = props.get('reg', 0.0) # Default 0
    cdef double betak = 1.0 # Scaling factor for xk.
    
    # Tracks for each entry of x, what iteration it was last updated at.
    cdef unsigned long[:] lag = np.zeros(m, dtype='I')
    
    # Initialize gradients
    cdef double gd = -1.0/n
    for i in range(n):
        indstart = indptr[i]
        indend = indptr[i+1]
        ydata = data[indstart:indend]
        yindices = indices[indstart:indend]
        ylen = indend-indstart
        add_weighted(gk, ydata, yindices, ylen, gd*b[i])

    # This is just a table of the sum the geometric series (1-reg/eta)
    # It is used to correctly do the just-in-time updating when
    # L2 regularisation is used.
    cdef double[:] lag_scaling = np.zeros(n*maxiter+1)
    lag_scaling[0] = 0.0
    lag_scaling[1] = 1.0
    cdef double geosum = 1.0
    cdef double mult = 1.0 - reg/eta
    for i in range(2,n*maxiter+1):
        geosum *= mult
        lag_scaling[i] = lag_scaling[i-1] + geosum
    
    # For least-squares, we only need to store a single
    # double for each data point, rather than a full gradient vector.
    # The value stored is the A_i * betak * x product
    cdef double[:] c = np.zeros(n)
    
    cdef unsigned long k = 0 # Current iteration number
    
    for epoch in range(maxiter):
            
        for j in range(n):
            if epoch == 0:
                i = j 
            else:
                i = np.random.randint(0, n)
            
            # Selects the (sparse) column of the data matrix containing datapoint i.
            indstart = indptr[i]
            indend = indptr[i+1]
            ydata = data[indstart:indend]
            yindices = indices[indstart:indend]
            ylen = indend-indstart
            
            # Apply the missed updates to xk just-in-time
            lagged_update(k, xk, gk, lag, yindices, ylen, lag_scaling, -1.0/(eta*betak))
            
            Aix = betak * spdot(xk, ydata, yindices, ylen)
            
            cnew = Aix
            cchange = cnew-c[i]
            c[i] = cnew
            betak *= 1.0 - reg/eta
            
            # Update xk with sparse step bit (with betak scaling)
            add_weighted(xk, ydata, yindices, ylen, -cchange/(eta*betak))
            
            k += 1
            
            # Perform the gradient-average part of the step
            lagged_update(k, xk, gk, lag, yindices, ylen, lag_scaling, -1.0/(eta*betak))
            
            # update the gradient average
            add_weighted(gk, ydata, yindices, ylen, cchange/n) 
    
    # Perform the just in time updates for the whole xk vector, so that all entries are up-to-date.
    gscaling = -1.0/(eta*betak)
    for ind in range(m):
        lagged_amount = k-lag[ind]
        lag[ind] = k
        xk[ind] += lag_scaling[lagged_amount]*gscaling*gk[ind]
    return betak * np.asarray(xk)
\end{python}



\end{document}
